% SPDX-FileCopyrightText: 2026 Yiyu Zhou <yiyu@yiyuzhou.io>
%
% SPDX-License-Identifier: GPL-3.0-or-later

\chapter{Saving and Show mode}
\label{ch:saving}

\section{Saving}

While a show has unsaved changes, the status bar shows \texttt{(unsaved)}
after its name, along with the key that saves it, \key{Control+s}.
\key{Control+S} saves the show to a different file, which the show uses from
then on, and \key{Control+o} opens another show.

cinq never discards unsaved changes:

\begin{enumerate}
\item In your own show, change the level of a cue, then close the main
  window without saving.
\item Start cinq again, press \key{Control+o}, and open
  \texttt{my-show.cinq.toml}.  The change is still there, and the status bar
  shows \texttt{(unsaved)}.
\item Press the undo key that the status bar names.  The show returns to its
  saved state.
\end{enumerate}

cinq records unsaved changes as they happen, in a separate file, so they
survive closing the window, opening another show, or a crash.  When the show
opens again, they return as a single change, which undo reverts to the saved
show.

\section{Show files and cue lists}

A show file stores the show as TOML, a plain-text format that is easy to
read, compare, and keep under version control.  Relative file paths in a show
are resolved from the show file's directory, so a show and its media can move
together.

A show can hold several cue lists, such as the acts of a play or a set of
sound effects to fire in any order:

\begin{enumerate}
\item In the storm show, click \ui{Effects} in the sidebar.  Its cues, E1
  and E2, are effects for the operator to fire as needed.
\item Select E2, Knock, press \ui{Standby selected}, and press \ui{GO}\@.
\item Click \ui{Act 1}.  Its standby has not moved.
\end{enumerate}

Each list has its own standby cue, and GO fires the standby cue of the list
shown.  Lists are added to a show file in a text editor, as the manual page
\texttt{cinq-show(5)} describes along with every other setting.

\section{Show mode}

Lock editing before a performance:

\begin{enumerate}
\item Check \ui{Show mode} in the transport.  The status bar shows
  \ui{Show mode} instead of \ui{Edit mode}, and the inspector becomes
  read-only.
\item Run the show as usual.  \ui{Add cue} and \ui{Delete cue} are disabled,
  and undo and redo only move standby back and forth.
\item Clear \ui{Show mode} to edit again.
\end{enumerate}

\section*{Exercise}

You can now protect a show: save it to a file, rely on cinq to keep unsaved
changes, and lock it for the performance.  Run the whole storm show in Show
mode, from cue~1 of Act~1 to the film in Act~2.
